\documentclass[a4paper,12pt]{article}
\usepackage[spanish,activeacute]{babel}
\usepackage[latin1]{inputenc}

%% Para las columnas
\usepackage{multicol}


\usepackage{amssymb, amsmath}
\usepackage{latexsym}
\usepackage{multicol}


%% Para numeraci?n autom?tica
\newcounter{etiqueta}
\newcounter{nroproblema}

\setcounter{nroproblema}{1}
\newcounter{nroapartado}
\setcounter{nroapartado}{1}
\newcommand{\n}{%
    \vspace{.1in}\noindent{\bf \arabic{nroproblema}.\ }%
    \addtocounter{nroproblema}{1}%
    \setcounter{nroapartado}{1}%
    }%
\newcommand{\apartado}{\vspace*{0.2cm}\hspace*{0.2cm}%
   {\bf \alph{nroapartado})}
   \addtocounter{nroapartado}{1}%
   }

%% m?genes

\setlength{\unitlength}{1mm} \addtolength{\voffset}{-15mm}
\addtolength{\textheight}{30mm} \addtolength{\hoffset}{-15mm}
\addtolength{\textwidth}{30mm}


\begin{document}

\begin{center}{\bf {\Large \'ALGEBRA I. HOJA 5}}\end{center}

\

Problemas de LADW: pp. 29-30, 6.1-6.12, p. 31, 7.1-7.5.

\


\n Estudiar los rangos de las siguientes matrices como funci�n de $\lambda$: 
$$\left(\begin{matrix}
2 &3 & 1& 7\\ 
3&  7 & -6 &-2\\ 
5& 8& 1 & \lambda
\end{matrix}\right)
, \left(\begin{matrix}
3& 1& 1& 4\\ 
\lambda &4& 10& 1\\ 
1& 7& 17& 3\\ 
2& 2& 4& 3 \end{matrix}\right)
, \left(\begin{matrix}
1 & \lambda & -1& 2\\ 
2& -1 &\lambda&  5\\ 
1 &10& -6& 1 \end{matrix}\right)$$



\n Estudiar si son subespacios vectoriales de $\mathbb{R}^2$ los siguientes subconjuntos: 

\apartado $S_1=\{(x, y) \in \mathbb{R}^2,\, y -x = 0\}$, 

\apartado $S_2 =\{(x, y) \in \mathbb{R}^2,\, y- x = 1\}$, 

\apartado $S_3 =\{(x, y) \in \mathbb{R}^2,\, xy = 0\}$, 

\apartado $S_4 =\{(x, y) \in \mathbb{R}^2,\, y = |x|\}$, 

\apartado $S_5 =\{(x, y) \in \mathbb{R}^2,\, |y| = |x|\}$,

\apartado $S_6=\{(x,y) \in \mathbb{R}^2,\, x\geq 0 \}$.

\vspace{0.2cm}

\n 
\apartado Comprobar que aunque $\mathbb{Q}$ est\'a contenido en $\mathbb{R}$ 
y sus operaciones est\'an inducidas por las de $\mathbb{R}$, $\mathbb{Q}$ no es subespacio vectorial de $\mathbb{R}$ 
(considerado como espacio vectorial de dimensi�n 1 sobre $\mathbb{R}$). 

\apartado Comprobar que $\mathbb{R}$ no es espacio vectorial sobre $\mathbb{C}$. 

\apartado Comprobar que $\mathbb{C}$ es un espacio vectorial de dimensi�n 2 sobre $\mathbb{R}$.


\vspace{0.2cm}

\n Estudiar si son subespacios vectoriales de $\mathbb{C}$, considerado como espacio vectorial complejo, 

\apartado el conjunto de los n\'umeros reales,

\apartado el conjunto de los n\'umeros imaginarios puros. 

\vspace{0.2cm}




\n Estudiar si son subespacios vectoriales del espacio vectorial real de funciones reales 
de variable real los siguientes subconjuntos: 

\apartado $S_1=\{f : f (0) = 0\}$,

\apartado $S_2 = \{f : f (1) = 0\}$,

\apartado $S_3 = \{f : f (0) = 1\}$. 

\vspace{0.2cm}

\n Averiguar si son subespacios vectoriales de los correspondientes espacios vectoriales los 
siguientes subconjuntos: 

\apartado  El subconjunto de los polinomios de variable real y con coeficientes reales, dentro de las 
funciones reales de variable real. 

\apartado El subconjunto de los polinomios de grado $\leq n$, de variable real y con coeficientes reales, 
siendo $n$ fijo, dentro del espacio vectorial de las funciones reales de variable real. 



\apartado El subconjunto de polinomios divisibles por $x-1$,  grado $\leq 3$,  de variable real y con 
coeficientes reales,  como subconjunto del espacio vectorial  de los polinomios de variable real, 
 grado $\leq 3$, y con coeficientes reales. 


\n Averiguar si son subespacios vectoriales de $\mathcal{M}_{2\times 2}(\mathbb{R})$ los siguientes subconjuntos: 

\apartado Las matrices $2\times 2$ con n\'umeros racionales.

\apartado Las matrices de n\'umeros reales de orden $2\times 2$ de traza cero. (Se llama traza de una 
matriz a la suma de los elementos de su diagonal). 

\apartado Las matrices de n\'umeros reales de orden $2\times 2$ de rango $1$. 

\apartado Las matrices de n\'umeros reales de orden $2\times 2$ que conmutan con la matriz 
B, siendo B una matriz fija $2\times 2$. 

\vspace{0.2cm}

\n Averiguar si los siguientes subconjuntos son subespacios vectoriales de $\mathcal{M}_{n\times n}(K)$: 

\apartado Las matrices diagonales.

\apartado Las matrices triangulares superiores. 

\apartado Las matrices sim\'etricas (las que satisfacen que $A^t=A$). 

\apartado Las matrices antisim\'etricas (las que satisfacen que $A=-A^t$).

\vspace{0.2cm}

\n Averiguar si el subconjunto de las matrices escalonadas de m filas y n columnas, 
es un subespacio vectorial del espacio  de matrices con entradas   reales (o complejas) de orden $m\times n$. 

\vspace{0.2cm}



\end{document}

\n Estudiar, seg\'un los valores de $\lambda$, la compatibilidad, o consistencia, de los sistemas $AX = b$, donde la matriz $(A|b)$ es:
$$\left(\begin{array}{ccc|c}
3& 2 &5 &1 \\
2 &4 &6& 3\\ 
5 &7& \lambda & 5 
\end{array}\right), \left(\begin{array}{ccc|c}
\lambda & 1 &  1& 1\\
 1 & \lambda & 1 & 1 \\
1& 1& \lambda & 1 \\
\end{array}\right).$$

\vspace{0.2cm}